\chapter{Analiza problemu i przegląd literatury}  % Analiza tematu

\note[hide]{
    To jest główny wstęp teoretyczny / state of the art. Wymagania mówią wprost, że analiza tematu ma zawierać wprowadzenie do dziedziny, sformułowanie problemu, poszerzone studia literaturowe, przegląd rozwiązań i algorytmów.
}

Zagadnienie rozpoznawania obiektów przez osoby niewidome wymaga uwzględnienia nie tylko aspektów technicznych, ale też społecznych i psychologicznych -- tego, jak użytkownik poznaje otoczenie i jakie informacje z niego wyciąga. Osoby widzące zazwyczaj nie mają problemów z rozpoznawaniem kształtów, kolorów, orientacji czy położenia obiektów w przestrzeni, natomiast osoby niewidome i słabowidzące w większym stopniu wykorzystują inne zmysły i technologie wspomagające, aby móc zbudować mentalną reprezentację otoczenia.

Ze względu na to projektowanie i wykonanie systemu rozpoznawania obiektów nie może skupiać się wyłącznie na aspektach technicznych, ale należy również uwzględnić perspektywę użytkownika docelowego. Należy zrozumieć, w jaki sposób informacja przekazywana przez komputer może wesprzeć proces poznawania otoczenia. W niniejszym rozdziale przedstawiono perspektywę osób niewidomych w kontekście rozpoznawania obiektów przestrzennych, a następnie omówiono wybrane technologie wspomagające osoby z niepełnosprawnością wzroku i metody komputerowej detekcji obiektów z wykorzystaniem zróżnicowanych danych wejściowych, takich jak obrazy RGB, mapy głębi czy dane LiDAR. Szczególną uwagę poświęcono rodzinie modeli YOLO, wykorzystywanej w niniejszej pracy. W kolejnej części rozdziału przedstawiono metryki wykorzystywane do oceny jakości detekcji obiektów, a następnie podsumowano przeprowadzoną analizę.

\section{Perspektywa osób niewidomych w rozpoznawaniu obiektów przestrzennych}

Przy projektowaniu systemu przeznaczonego dla osób niewidomych należy wziąć pod uwagę różnice w percepcji przestrzennej pomiędzy osobami widzącymi a niewidomymi. W literaturze zagadnienie to jest analizowane zarówno z perspektywy technicznej, psychologicznej, jak i pedagogicznej. Badania nad percepcją osób niewidomych pozwalają lepiej zrozumieć ich sposób poznawania otoczenia. Costa et al. w swojej pracy, na podstawie studium przypadku, pokazują, że rozpoznawanie obiektu jest całym procesem stopniowego pozyskiwania i integrowania informacji. Zaznaczają, że w przypadku osób widzących spojrzenie zazwyczaj pozwala zrozumieć obiekt całościowo, natomiast w przypadku osób niewidomych poznanie odbywa się sekwencyjnie, etapami. Duże znaczenie mają więc cechy obiektu, które są wyraźnie dostępne w eksploracji dotykowej: prostota, regularność, krawędzie oraz faktura. Wspierają one budowę mentalnej reprezentacji przestrzennej obiektu. Osoba niewidoma może opierać ją nie tylko na aktualnej informacji dotykową, ale również wcześniejszych doświadczeniach, kontekście i wiedzy o świecie. Autorzy ukazują, że brak wzroku nie oznacza prostego zastąpienia go przez inne zmysły -- dotyk, słuch, pamięć i doświadczenia dostarczają informacji innego rodzaju oraz są wykorzystywane w odmienny sposób \cite{bib:spacialPerception}.

Rozwinięcie tej myśli znajdujemy w pracy autorstwa Figueiras i Arcavi, którzy pokazują jak osoba niewidoma uczy się matematyki i w jaki sposób komunikuje się z nauczycielem. Autorzy opisują przebieg lekcji prowadzonych przez niewidomego nauczyciela dla trzech uczniów, z czego dwóch było w pełni niewidomych, a jeden miał częściową utratę wzroku. Szczególnie istotne okazują się zapisy rozmów pomiędzy nauczycielem a uczniami, które ukazują wcześniej wspomniane etapowe poznawanie obiektu. Podczas omawiania stożka nauczyciel bazuje na wcześniejszych doświadczeniach i wiedzy uczniów, a także na ich wyobrażeniach przestrzennych.

W kolejnej części artykułu autorzy ukazują w jaki sposób nauczyciel prowadzi uczniów przez analizę wykresu funkcji kwadratowej. Nauczyciel celowo odracza rozwiązanie algebraiczne i prowadzi uczniów przez proces rozumowania. Początkowo uczestnicy rozważają położenie jednego z pierwiastków, następnie wierzchołka i osi symetrii. Dopiero po tym, gdy uczniowie przeprowadzili rozumowanie, nauczyciel pozwala na rozwiązanie algebraiczne. Pokazuje to, że w przypadku, gdy nie mamy dostępu do globalnej reprezentacji obrazu, uwypukla się znaczenie jawnych opisów relacji pomiędzy charakterystycznymi cechami obiektu, odwołania do wcześniejszych reprezentacji i uporządkowanie procesu rozumowania \cite{bib:learningToSee}. W kontekście projektowanego systemu sugeruje to, że sama informacja o klasie rozpoznanego obiektu nie jest wystarczająca, a zasadne może być również przekazywanie informacji o jego cechach charakterystycznych i zachodzących między nimi relacjach.

Odchodząc od ściśle teoretycznych i obserwacyjnych badań, w literaturze znajdujemy również przykłady artykułów naukowych, które bezpośrednio wskazują problemy osób niewidomych. Güler i Ürey przeprowadzili badanie jakościowe obejmujące 13 uczniów szkoły dla osób z niepełnosprawnością wzroku oraz ich 10 nauczycieli. Dotyczyło ono wyzwań i potrzeb uczniów przy nauce przedmiotów przyrodniczych. Istotnym wnioskiem z dokumentu jest wskazanie, że nauczyciele częściej korzystali z podręczników pisanych alfabetem Braille'a oraz tabliczek i rylców do zapisu brajlowskiego, aniżeli z materiałów przestrzennych, modeli 3D czy eksperymentów. Jeden z uczniów wskazywał, że ma problem z wyobrażaniem sobie reprezentacji przestrzennych słów, twierdząc, że gdy słyszy np. słowo "piramida" nie potrafi zbudować w umyśle odpowiadającego mu obrazu. Nauczyciele wskazują również, że tureccy uczniowie dotknięci niepełnosprawnością wzroku nie mają przewidzianego osobnego toku nauczania. Są więc zmuszeni do konkurowania z widzącymi rówieśnikami tak, jakby byli w stanie wyuczyć się wszystkiego w ten sam sposób. Zaznaczają też, że zagadnienia związane z naukami przyrodniczymi są często abstrakcyjne, co utrudnia ich zrozumienie, a sam proces nauki powinien być dostosowany do potrzeb uczniów z niepełnosprawnością wzroku \cite{bib:awareNeeds}. Podobne problemy związane z brakiem materiałów edukacyjnych, koniecznością dostosowania istniejących, wydłużonym czasem przygotowania zajęć i dostosowania komunikacji odnotowano również w badaniach dotyczących nauki języka angielskiego \cite{bib:studentTeacherPerception}.

Kolejnym zagadnieniem dotyczącym perspektywy osób niewidomych, na które zwrócili uwagę Li et al., jest sposób, w jaki różne źródła informacji: dotyk, opis dźwiękowy, wcześniejsze doświadczenia; wspierają rozumienie obiektów. Zaznaczają, że osoby niewidome nie bazują jedynie na pojedynczym źródle informacji, ale łączą informacje pochodzące z różnych modalności. Badana grupa osób wskazała, że poprzez dotyk mogą zrozumieć kształt obiektu, jego położenie, formę i szczegóły przestrzenne, ale by w pełni zrozumieć kontekst, historię lub znaczenie obiektu, potrzebują dodatkowego opisu audio. Dodatkowo, autorzy zwracają uwagę na to, że osoby niewidzące różnią się między sobą pod względem preferencji i sposobów poznawania obiektów. Zaznaczają, że różnią się one w zależności od stopnia niepełnosprawności wzroku oraz tego, czy osoba jest niewidoma od urodzenia, czy utraciła wzrok w późniejszym okresie życia \cite{bib:blindVisualArts}. Na tej podstawie można przyjąć, że projektowany system nie powinien ograniczać się do jednego sposobu komunikacji, ale pozwolić na personalizację i dostosowanie do preferencji konkretnego użytkownika

Przytoczone badania wskazują, że rozpoznawanie obiektów przez osoby niewidome nie jest pojedynczym aktem identyfikacji, ale etapowym procesem poznawczym. Eksploracja dotykowa pozwala na poznanie lokalnych części obiektu: krawędzi, powierzchni, faktury i orientacji, ale to wcześniejsze doświadczenia i przekaz słowny pozwalają na zbudowanie jego pełnej reprezentacji mentalnej. Jednocześnie, mała ilość dedykowanych materiałów edukacyjnych mocno utrudnia procesy nauczania -- zwłaszcza przy pojęciach abstrakcyjnych.

Wynika z tego, że system wspomagający nie powinien ograniczać się jedynie do przekazywania nazw obiektów, ale wspierać użytkownika w całym procesie jego rozpoznawania. Informacja zwrotna powinna być jednoznaczna, możliwa do dostosowania do potrzeb użytkownika i wspierać go w naturalnym procesie eksploracji dotykowej. Te wnioski są podstawą do analizy istniejących rozwiązań technologicznych przedstawionej w kolejnym podrozdziale.

\section{Technologie wspomagające osoby z niepełnosprawnością wzroku}

\note[hide]{
    Aplikacje mobilne, systemy wearable, kamery, czujniki głębi, audio feedback.
}

Analiza perspektywy osób niewidomych pozwala zrozumieć, że wsparcie dla nich

\subsection{Rozwój technologii asystujących}

Naukowcy od dawna widzieli w rozwoju techniki okazję do wsparcia osób z jakąkolwiek niepełnosprawnością. Jednym z pierwszych urządzeń wspomagających niesłyszących był aparat słuchowy opracowywany przez Millera Reese'a Hutchisona od około 1895 roku \cite{bib:hearingAidsHistory}. Rozwiązaniem, które zapowiadało późniejszy rozwój elektronicznych systemów wspomagających był system \nobreakdash{POSSUM} (ang. \textit{Patient Operated Selector Mechanism}) opracowany początkiem lat 60. XX wieku. Umożliwiał on osobom z poważną niepełnosprawnością ruchu sterowanie maszyną do pisania \cite{bib:possum}. Nie można jednak mówić o cyfrowych systemach wspomagających przed pojawieniem się elektronicznych komputerów. W latach 60. i 70. XX wieku zaczęły pojawiać się bardziej zaawansowane systemy wspomagające oparte o cyfrowe jednostki obliczeniowe. Przykładem jednego z nich może zostać Tufts Interactive Communicator, który umożliwiał niemym z ograniczeniami ruchowymi komunikację z otoczeniem poprzez wybieranie znaków za pomocą mechanizmu skanowania, a następnie prezentację pełnego komunikatu na wyświetlaczu w formie wydruku \cite{bib:augmentativeCommunication}. Późniejszy przykład podobnego systemu był wykorzystywany przez wybitnego fizyka Stephena Hawkinga który, pomimo swojej postępującej niepełnosprawności ruchowej, mógł sprawnie komunikować się z otoczeniem poprzez niewielkie ruchy policzka \cite{bib:intel2014hawking}.

Równocześnie rozwijały się technologie wspomagające osoby z niepełnosprawnością wzroku, w którym dane przekształcano na formę możliwą do odbioru przez inne zmysły. Kamieniem milowym okazało się urządzenie \nobreakdash{Optacon} (ang. \textit{Optical to Tactile Converter}) opatentowane w 1966 roku przez Johna G. Linvilla. Umożliwiało ono przekształcenie obrazu na wibracje matrycy, które mogły być odczytane przez użytkownika za pomocą dotyku \cite{pat:Linvill1966ReadingAid}. W kolejnych dekadach rozwijano m. in. systemy rozpoznawania znaków, syntezatory mowy, czytniki ekranów i urządzenia mobilne, stopniowo zwiększając i umożliwiając osobom niewidomym dostęp do informacji cyfrowych.

Współczesny rozwój takich rozwiązań dokonuje się dzięki stopniowej miniaturyzacji elektroniki i popularyzacji urządzeń osobistych. Według prognozy MarketsandMarkets rynek urządzeń ubieralnych ma wzrosnąć z 87,19 mld dolarów w 2025 roku do 238,71 mld dolarów w 2032 przy wzroście średniorocznym o 14,7\% \cite{bib:marketsandmarkets}. Dane te są przesłanką do wniosku, że w najbliższych latach możemy spodziewać się dynamicznego rozwoju małych urządzeń wyposażonych w czujniki, układy umożliwiające komunikację i mikroprocesory zdolne do szybkiego, lokalnego przetwarzania danych. Te same podzespoły mogą umożliwiać również rozwój systemów wspomagających osoby niewidome. Dedykowane urządzenia, noszone na ciele użytkownika, będą mogły stale analizować otoczenie, a następnie przekazywać użytkownikowi informacje poprzez komunikaty dźwiękowe czy haptyczne. Dodatkowo wykorzystanie szybkich sieci bezprzewodowych i nagły rozwój uczenia maszynowego pozwoli na szybsze tworzenie i udoskonalanie działających w czasie rzeczywistym algorytmów.

\subsection{Standardy i regulacje prawne dotyczące dostępności}

Rozwój technologiczny nie był jedynym czynnikiem, który pchnął do przodu upowszechnienie się systemów wspomagających. Istotną rolę odegrały również regulacje prawne i standardy, które nakładały obowiązek zapewnienia dostępności systemów informatycznych dla osób z niepełnosprawnością. Bardzo istotnym krokiem było opublikowanie w 1999 roku przez World Wide Web Consortium (W3C) pierwszej wersji wytycznych WCAG (ang. \textit{Web Content Accessibility Guidelines}) określających, w jaki sposób powinny być projektowane treści internetowe, by były one dostępne dla osób z różnymi rodzajami niepełnosprawności \cite{bib:wcag1}. Kolejne wersje wytycznych uporządkowały zasady wokół czterech głównych: postrzegalności, funkcjonalności, zrozumiałości i solidności. Najnowsza wersja standardu WCAG~2.2 została opublikowana 12 grudnia 2024 \cite{bib:wcag22}. Obecnie trwają prace nad kolejną wersją WCAG~3.0.

Znaczenie dostępności zostało następnie zwiększone dzięki wdrożeniu kroków prawnych. Artykuł 9 Konwencji ONZ o prawach osób z niepełnosprawnościami nakłada na państwa członkowskie obowiązek zapewnienia dostępu do informacji, komunikacji, infrastruktury i technologii do potrzeb osób z niepełnosprawnością \cite{bib:crpd}. Unia Europejska wprowadziła w 2016 roku dyrektywę 2016/2102, która określa wymagania dotyczące dostępności stron internetowych i aplikacji mobilnych podmiotów sektora publicznego \cite{bib:directive2016accessibility}. Polskie prawo wprowadziło te rozporządzenie w życie ustawą z dnia 4 kwietnia 2019 roku o dostępności cyfrowej stron internetowych i aplikacji mobilnych podmiotów publicznych \cite{bib:ustawaDostepnosc2019}. Od 28 czerwca 2025 roku w Polsce obowiązuje również Europejski Akt o Dostępności, który nakłada wymagania na wybrane usługi i produkty, takie jak urządzenia mobilne, komputery, czytniki książek, handel elektroniczny, bankowość i transport \cite{bib:eea2019, bib:polskiAktDostepnosci2024}. Szczegółowe wymagania techniczne dla produktów i usług ICT określa norma EN~301~549, która jest stosowana w całej Unii Europejskiej \cite{bib:en301549}.

Dla osób niewidomych te regulacje mają szczególne znaczenie w kontekście współpracy czytników ekranowych z interfejsami, możliwości użycia bez użycia wzroku i jednoznacznego przekazywania informacji. Rozwój prawa sprawił więc, że dostępność cyfrowa nie pozostała w fazie teoretycznej, ale stała się wymogiem, a przez to istotnym czynnikiem technologii wspomagających.

\subsection{Współczesne systemy wspomagające osoby niewidome}

\note{
    \begin{enumerate}
        \item Potdar, Pai i Akolkar — A Convolutional Neural Network based Live Object Recognition System as Blind Aid \cite{bib:cnnBlindAid}\\
              Dałbym go jako prosty punkt wyjścia. System wykorzystuje kamerę i sieć CNN do rozpoznawania obiektów w czasie rzeczywistym. Pokazuje klasyczne podejście:

              obraz z kamery → rozpoznanie obiektu → przekazanie informacji użytkownikowi.

              Artykuł jest starszy i technicznie prostszy od późniejszych rozwiązań, dzięki czemu nadaje się do pokazania, jak rozwijały się komputerowe technologie wspomagające. Nie trzeba szczegółowo opisywać ImageNet ani architektury Inception — to zostawiamy do przeglądu metod.
        \item Sethuraman et al. — MagicEye: An Intelligent Wearable Towards Independent Living of Visually Impaired \cite{bib:magiceye}\\
              To jeden z najbardziej bezpośrednio związanych z Twoją pracą artykułów. Przedstawia urządzenie ubieralne z kamerą, które wykorzystuje model oparty na YOLOv5 do rozpoznawania przedmiotów codziennego użytku w czasie rzeczywistym.

              Jego rola w 2.2:

              pokazanie przejścia od eksperymentalnego modelu do kompletnego systemu;
              mobilność i urządzenie znajdujące się przy użytkowniku;
              wymaganie działania w czasie rzeczywistym;
              ograniczenia związane z liczbą rozpoznawanych klas i jakością obrazu.

              To dobre bezpośrednie odniesienie dla Twojego układu: urządzenie mobilne, obraz z kamery, model detekcji i informacja dla użytkownika.
        \item Baig et al. — AI-based Wearable Vision Assistance System for the Visually Impaired \cite{bib:wearablevisionassistance}\\
              Ten artykuł pokazuje kolejny etap rozwoju: system nie ogranicza się do samej etykiety obiektu. Łączy detekcję MobileNet-SSD, rozpoznawanie twarzy oraz model vision-language generujący kontekstowy opis sceny.

              Bardzo dobrze łączy się to z wnioskami z 2.1:

              sama informacja „to jest obiekt X” może być niewystarczająca; użytkownik może potrzebować również informacji o kontekście, cechach i otoczeniu.

              Warto opisać też:

              działanie na urządzeniu Raspberry Pi;
              średni czas odpowiedzi;
              pogorszenie skuteczności przy słabym oświetleniu;
              możliwość dodawania nowych obiektów i osób.

              Ponieważ jest to preprint, przedstawiałbym go jako przykład prototypowego rozwiązania, a nie jako ostateczne potwierdzenie skuteczności całej klasy systemów.
        \item Wang et al. — YOLO-OD: Obstacle Detection for Visually Impaired Navigation Assistance \cite{bib:yolo_od}\\
              Ten artykuł reprezentuje drugi duży obszar technologii wspomagających: nawigację i wykrywanie przeszkód. Nie dotyczy bezpośrednio rozpoznawania dotykanych brył, ale pokazuje, że detekcja obiektów jest wykorzystywana nie tylko do ich identyfikacji, lecz także do wspierania orientacji przestrzennej.

              Wystarczy opisać:

              wykrywanie przeszkód w scenach zewnętrznych;
              zastosowanie zmodyfikowanego modelu YOLO;
              konieczność wykrywania obiektów o różnych rozmiarach;
              znaczenie kompromisu między skutecznością a szybkością działania.

              Nie wchodziłbym tutaj głęboko w bloki FWB, ABB czy EFAH. To podrozdział o technologiach wspomagających, nie o szczegółach architektury.
        \item Dzierzgowska et al. — Alternative audio-graphic method for presenting structural information in mathematical graphs designed for low-vision users \cite{bib:Dzierzgowska2025}\\
              Ten artykuł nie opisuje klasycznej detekcji obiektów, ale jest bardzo ważny ze względu na sposób przekazywania informacji użytkownikowi. Pokazuje rozwiązanie łączące grafikę z informacją dźwiękową i zawiera badanie użytkowników z wykorzystaniem czasu wykonania zadań, liczby błędów, NASA-TLX i SUS.

              Jego rola:

              pokazanie, że technologia wspomagająca to nie tylko model rozpoznający obiekt;
              znaczenie dobrze zaprojektowanej informacji zwrotnej;
              ocena obciążenia, frustracji i użyteczności;
              uzasadnienie wykorzystania komunikatów dźwiękowych w Twoim systemie.

              To źródło jest mocne, bo przedstawia nie tylko rozwiązanie techniczne, ale też jego ocenę z udziałem użytkowników.
    \end{enumerate}
    Dodatkowo: Kawulok i Maćkowski — YOLO-Type Neural Networks in the Process of Adapting Mathematical Graphs to the Needs of the Blind \cite{bib:yolo_adaptation}\\
    Pokazuje zastosowanie modeli YOLO do automatycznej adaptacji materiałów graficznych dla osób niewidomych. Dobrze pasuje jako przykład technologii edukacyjnej i jako łącznik między 2.1 a techniczną częścią pracy.

    Nie poświęcałbym mu jednak całego dużego akapitu, ponieważ szczegółowe porównanie modeli YOLO będzie lepiej pasowało do 2.5.
}

\draft{
    WNIOSEK KOŃCOWY:

    Współczesne technologie wspomagające coraz częściej łączą mobilną akwizycję obrazu, detekcję działającą w czasie rzeczywistym oraz dźwiękową lub kontekstową informację zwrotną. Jednocześnie większość opisanych systemów koncentruje się na rozpoznawaniu obiektów znajdujących się w otoczeniu lub na wykrywaniu przeszkód. Mniej uwagi poświęca się rozpoznawaniu obiektu, który jest aktualnie eksplorowany dotykowo przez użytkownika.
}

\section{Problem rozpoznawania obiektów 3D w interakcji dotykowej}

\note{
    Tutaj opisuję, że użytkownik nie tylko patrzy kamerą na scenę, ale dotyka obiektu, więc problem jest trochę inny niż klasyczna detekcja przedmiotów w pomieszczeniu.
}

\section{Metody detekcji obiektów w obrazach}

\note{
    YOLO, SSD, Faster R-CNN, Mask R-CNN — ogólnie, bez przesadnego grzebania.
}

\section{Rodzina modeli YOLO}

\note{
    Tu najwięcej, bo tego używasz. Architektura ogólnie, zalety: szybkość, detekcja w czasie rzeczywistym, różne rozmiary modeli.
}

\section{Wykorzystanie danych RGB, głębi i LiDAR w rozpoznawaniu obiektów}

\note{
    Tu mogę opisać, dlaczego LiDAR był rozważany, ale w praktyce ma ograniczenia przy małych obiektach i małych odległościach.
}

\section{Metryki oceny detekcji obiektów}

Ocena modeli detekcji obiektów wymaga zastosowania odpowiednich metryk, które pozwalają na porównanie uzyskanych wyników. Metryki te powinny uwzględniać zarówno poprawność klasyfikacji, jak i dokładność lokalizacji obiektu na obrazie. Przeciwnie do klasyfikacji obrazów, przy detekcji nie wystarczy, aby model poprawnie rozpoznał klasę obiektu. Równie istotne jest, aby poprawnie wyznaczył ramkę ograniczającą wokół wykrytego obiektu. Jedną z najczęściej stosowanych metryk jest IoU (ang. \textit{Intersection over Union}), określająca stopień pokrycia między ramką obiektu przewidywaną przez model a ramką rzeczywistą \cite{bib:detectionMetrics}. Wartość IoU obliczana jest jako stosunek pola części wspólnej obu ramek do pola ich części łącznej:

\begin{equation}
    IoU = \frac{|B_p \cap B_{gt}|}{|B_p \cup B_{gt}|},
\end{equation}

gdzie $B_p$ jest ramką przewidywaną przez model, a $B_{gt}$ jest ramką referencyjną (ang. \textit{ground truth}). Im większa wartość IoU, tym lepsze dopasowanie ramki przewidywanej do ramki rzeczywistej. W praktyce najczęściej stosuje się progi, np. $IoU \geq 0.5$, aby zdecydować, czy dana detekcja może być uznana za poprawną. Próg ten został wykorzystany w wielu klasycznych procedurach oceny detekcji, takich jak PASCAL VOC \cite{bib:PASCALVOCChallenge}, i jest nadal powszechnie stosowany w wielu benchmarkach. Na podstawie wartości IoU oraz zgodności klasy obiektu możliwe jest przypisanie wyników detekcji do jednej z kategorii: prawdziwych pozytywów (TP, ang. \textit{true positive}), fałszywych pozytywów (FP, ang. \textit{false positive}) i fałszywych negatywów (FN, ang. \textit{false negative}). Detekcja jest traktowana jako TP, jeżeli model poprawnie rozpoznał klasę obiektu, a przewidywana ramka spełnia zadany próg IoU względem ramki referencyjnej. FP to detekcja błędna w której model przewidział klasę obiektu, ale nie spełnił progu IoU lub przewidziana klasa była błędna. FN natomiast to sytuacja, w której model nie wykrył obiektu, który rzeczywiście był na obrazie.

Jedną z podstawowych metryk oceny jest precyzja (ang. \textit{precision}), określająca udział poprawnych detekcji modelu wśród wszystkich detekcji, jakie zwrócił:

\begin{equation}
    Precision = \frac{TP}{TP + FP}.
\end{equation}

Wysoka precyzja oznacza, że model generuje niewiele fałszywych pozytywów. Drugą, powiązaną z nią metryką, jest czułość (ang. \textit{recall}), określająca, jaki jest udział poprawnych detekcji wśród wszystkich rzeczywistych obiektów:

\begin{equation}
    Recall = \frac{TP}{TP + FN}.
\end{equation}

Obie metryki pozostają ze sobą w relacji kompromisu, ponieważ zwiększenie progu pewności detekcji może ograniczyć liczbę fałszywych pozytywów, ale jednocześnie prowadzić do pominięcia części rzeczywistych obiektów. Z kolei obniżenie tego progu może zwiększyć liczbę detekcji, ale równocześnie powodować wzrost liczby fałszywych detekcji. Z tego względu, przy analizie detektorów często łączy się je w jedną metrykę -- krzywą PR (precyzja-czułość, ang. \textit{precision-recall curve}), przedstawiającą zależność pomiędzy nimi dla różnych wartości progu decyzyjnego. Na jej podstawie wyznacza się metrykę AP (ang. \textit{Average Precision}), odpowiadającą polu pod krzywą PR.

W przypadku detekcji wieloklasowej używa się średniej wartości AP obliczonej dla wszystkich klas, określanej jako mAP (ang. \textit{mean Average Precision}):

\begin{equation}
    mAP = \frac{1}{N} \sum_{i=1}^{N} AP_i,
\end{equation}

gdzie $N$ jest liczbą klas, a $AP_i$ jest wartością AP dla $i$-tej klasy. W praktyce najczęściej wykorzystuje się dwie wersje tej metryki: $mAP50$, gdzie detekcję uznaje się za poprawną przy progu $IoU \geq 0.5$, oraz $mAP50-95$, gdzie wynik jest uśredniany dla progów $IoU$ z zakresu od $0.5$ do $0.95$. Pierwsza z nich jest bardziej liberalna i pozwala określić ogólną zdolność modelu do wykrycia obiektu. Druga natomiast jest bardziej rygorystyczna, ponieważ wymaga dokładniejszego dopasowania ramek ograniczających, dzięki czemu lepiej pokazuje zarówno skuteczność klasyfikacji, jak i lokalizacji obiektu. Uśrednianie wyników dla wielu progów zostało szeroko zastosowane m. in. w metodyce oceny detekcji zbioru Microsoft COCO \cite{bib:microsoftCOCO} i jest obecnie standardem w wielu benchmarkach detekcji obiektów.

Dodatkowo, w pracy wykorzystano metrykę $F1$-score, będącą średnią harmoniczną precyzji i czułości:

\begin{equation}
    F1 = 2 \cdot \frac{P \cdot R}{P + R},
\end{equation}

gdzie $P$ jest precyzją, a $R$ jest czułością. Metryka ta umożliwia przedstawienie kompromisu pomiędzy liczbą fałszywych detekcji a liczbą pominiętych obiektów. Może być szczególnie przydatna, gdy sama precyzja lub czułość nie dają pełnego obrazu jakości modelu. W niniejszej pracy metrykę $F1$-score wykorzystano jako metrykę pomocniczą względem podstawowych metryk $mAP50$ i $mAP50-95$, aby lepiej zrozumieć kompromis pomiędzy precyzją a czułością w kontekście detekcji obiektów dotykanych przez użytkownika.

Oprócz wartości liczbowych wykorzystano również macierze pomyłek (ang. \textit{confusion matrix}), pozwalające przeanalizować, które klasy obiektów są najczęściej mylone przez model. Są one przedstawiane w formie tabel, w których jedna oś odpowiada rzeczywistym klasom, a druga klasom przewidywanym przez model. Dla dobrze działającego modelu większość wartości powinna znajdować się na przekątnej macierzy, co wskazuje na poprawną klasyfikację obiektów. Macierze pomyłek są szczególnie istotne przy rozpoznawaniu brył geometrycznych o podobnych kształtach, np. sześcianu i prostopadłościanu, ponieważ sama wartość mAP nie wskazuje bezpośrednio, pomiędzy którymi klasami model popełnia najwięcej błędów. Analiza macierzy pomyłek może również pomóc w identyfikacji problemów z danymi treningowymi, np. niedostatecznej liczby przykładów dla wybranych klas lub błędów w oznaczeniach.

\section{Podsumowanie analizy tematu}

\note{
    Najważniejsze: z literatury wynika, że detekcja obiektów w czasie rzeczywistym jest sensownym kierunkiem, ale Twoja nisza to rozpoznawanie konkretnych dotykanych brył 3D w kontrolowanym scenariuszu.
}

% 

%\begin{itemize}
%\item Jaki problem chcę (muszę :-) rozwiązać?
%\item Dlaczego rozwiązanie problemu jest ważne?
%\item Jak inni rozwiązują ten problem?
%\item Jakie są zalety i wady tych rozwiązań?
%\end{itemize}

% Odwołania do literatury: 
% książek \cite{bib:ksiazka},
% artykułów w czasopismach \cite{bib:artykul}, 
% materiałów konferencyjnych \cite{bib:konferencja}
% i stron www \cite{bib:internet}.

% Równania powinny być numerowane
% \begin{align}
% y = \frac{\partial x}{\partial t}
% \end{align}

%
%\begin{itemize}
%\item analiza tematu
%\item wprowadzenie do dziedziny (\english{state of the art}) – sformułowanie problemu, 
%\item poszerzone studia literaturowe, przegląd literatury tematu (należy wskazać źródła wszystkich informacji zawartych w pracy)
%\item opis znanych rozwiązań, algorytmów, osadzenie pracy w kontekście
%\item Tytuł rozdziału jest często zbliżony do tematu pracy. 
%\item Rozdział jest wysycony cytowaniami do literatury \cite{bib:artykul,bib:ksiazka,bib:konferencja}. 
%Cytowanie książki \cite{bib:ksiazka}, artykułu w czasopiśmie \cite{bib:artykul}, artykułu konferencyjnego \cite{bib:konferencja} lub strony internetowej \cite{bib:internet}.
%\end{itemize}
